% ============================================================
\section{Funtor \texorpdfstring{$G: Par \to Rel$}{G: Par -> Rel}}
\label{sec:functor-par-rel}

\subsection{Descrição das Categorias}

\begin{definition}[Categoria $\mathbf{Par}$ (modelo MATLAB/Octave)]
Os objetos são conjuntos finitos representados pelo seu cardinal $n$. Um morfismo parcial $f: n \rightharpoonup m$ é um vetor de comprimento $n$ com entradas em $\{0, 1, \ldots, m\}$, onde $0$ indica que o elemento não tem imagem definida.
\end{definition}

\begin{lstlisting}[style=matlab, caption={Objeto e morfismo em $\mathbf{Par}$}]
n = 3;  m = 4;
% f(1)=2, f(2)=indefinido, f(3)=1
f = [2, 0, 1];
\end{lstlisting}

\begin{definition}[Categoria $\mathbf{Rel}$ (modelo MATLAB/Octave)]
Os objetos são conjuntos finitos. Um morfismo $R: n \to m$ é uma relação binária $R \subseteq \{1,\ldots,n\} \times \{1,\ldots,m\}$, representada por uma matriz booleana $m \times n$ onde a entrada $(i,j)$ vale $1$ se e só se $(j, i) \in R$.
\end{definition}

\begin{lstlisting}[style=matlab, caption={Objeto e morfismo em $\mathbf{Rel}$}]
% Relacao R ⊆ {1,2,3} x {1,2,3,4}: matriz booleana 4x3
M_rel = [0, 0, 1;   % f(3)=1
         1, 0, 0;   % f(1)=2
         0, 0, 0;
         0, 0, 0];
\end{lstlisting}

\subsection{Funtor Covariante \texorpdfstring{$G: Par \to Rel$}{G: Par -> Rel}}

O funtor $G$ transforma aplicações parciais nas suas relações gráficas associadas. Qualquer função parcial pode ser vista como uma relação binária cujo grafo omite os pares correspondentes a elementos sem imagem definida.

\begin{itemize}
    \item \textbf{Nos objetos:} $G(n) = n$. O conjunto mantém-se inalterado.
    \item \textbf{Nos morfismos:} Um morfismo parcial $f: n \rightharpoonup m$ é mapeado para a relação binária $R_f = \{(x,\, f(x)) \mid f(x) \neq 0\}$. Qualquer elemento indefinido não gera uma entrada com valor $1$ na coluna correspondente da matriz booleana.
\end{itemize}

\begin{lstlisting}[style=matlab, caption={Funtor $G$ aplicado a uma aplicação parcial $f: 3 \rightharpoonup 4$}]
f = [2, 0, 1];   % f(1)=2, f(2)=indefinido, f(3)=1
n = length(f);
m = max(f);      % cardinal do codominio

% Construcao da matriz relacional correspondente
M_rel = zeros(m, n);
for j = 1:n
    if f(j) ~= 0
        M_rel(f(j), j) = 1;
    end
end
% Resultado: M_rel = [0,0,1; 1,0,0; 0,0,0; 0,0,0]
\end{lstlisting}

\begin{proposition}
O funtor $G: \mathbf{Par} \to \mathbf{Rel}$ satisfaz as seguintes propriedades:
\begin{itemize}
    \item $G(\mathrm{id}_n)$ é a matriz identidade $n \times n$ (aplicação total que preserva todos os elementos);
    \item $G$ preserva a composição de aplicações parciais: $G(g \circ f) = G(g) \circ_{\mathbf{Rel}} G(f)$, onde a composição em $\mathbf{Rel}$ é a composição relacional;
    \item $G$ é fiel: aplicações parciais distintas produzem relações distintas.
\end{itemize}
\end{proposition}

\subsection{Transformações Naturais \texorpdfstring{$\beta: G_1 \Rightarrow G_2$}{beta: G1 => G2}}

Dados dois funtores paralelos $G_1, G_2: \mathbf{Par} \to \mathbf{Rel}$, uma transformação natural $\beta: G_1 \Rightarrow G_2$ é uma família de morfismos em $\mathbf{Rel}$,
\[
    \beta_n : G_1(n) \to G_2(n), \quad \text{para cada objeto } n \in \mathbf{Par},
\]
tal que para cada morfismo $f: n \rightharpoonup m$ o seguinte quadrado comuta:
\[
    G_2(f) \circ \beta_n \;=\; \beta_m \circ G_1(f).
\]
Considerando $G_1$ como o funtor que mapeia a aplicação parcial para o seu grafo exato (matriz esparsa) e $G_2$ como o funtor que mapeia para a relação universal (matriz totalmente preenchida a $1$), a transformação natural $\beta: G_1 \Rightarrow G_2$ verifica que o grafo exato é uma sub-relação da relação universal.

\begin{lstlisting}[style=matlab, caption={Verificação da transformação natural $\beta: G_1 \Rightarrow G_2$ em $\mathbf{Par} \to \mathbf{Rel}$}]
n = 3;  m = 4;
f = [2, 0, 1];   % aplicacao parcial

% Funtor G1: grafo exato da aplicacao parcial
M_G1 = zeros(m, n);
for x = 1:length(f)
    if f(x) ~= 0, M_G1(f(x), x) = 1; end
end

% Funtor G2: relacao universal
M_G2 = ones(m, n);

% Validacao: G1 e sub-relacao de G2 (condicao de naturalidade)
beta_n = all((M_G1 & M_G2) == M_G1, 'all');  % Retorna 1 (true)
\end{lstlisting}